\documentclass[a4paper,fontset=none]{ctexart}

\usepackage{indentfirst}
\setlength{\parindent}{0em}
\usepackage{autobreak}
\allowdisplaybreaks
\usepackage{parskip}
\usepackage{geometry}
\geometry{top=3cm,bottom=3cm,left=2cm,right=2cm}
\usepackage{anyfontsize}

\usepackage{fontspec}
\setmainfont{STIX Two Text}
\usepackage{unicode-math}
\unimathsetup{mathrm=sym,mathbf=sym,mathsf=sym,bold-style=ISO}
\setmathfont{STIX Two Math}
\usepackage{physics2}
\usephysicsmodule{ab,op.legacy}
\usepackage{fixdif,derivative}
\usepackage{tabularray}
\UseTblrLibrary{amsmath}
\newcommand{\um}{\upmu\mathrm{m}}
\newcommand{\ang}{\text{\r{A}}}
\AtBeginDocument{
    \let\uglyepsilon\epsilon
    \renewcommand\epsilon{\varepsilon}
}

\setCJKmainfont{Source Han Serif CN}[BoldFont=Source Han Sans CN Medium,ItalicFont=FZKai-Z03]

\title{\textbf{星际介质天文学作业}}
\author{张植竣}
\date{2024年5月31日}

\begin{document}

\maketitle

\begin{enumerate}
    \item[18.1]
    由(18.9)式：
    \[
        \frac{I([\mathrm{O\,III}]5008)}{I(\mathrm{H}\beta)} = \frac{n(\mathrm{O\,III})}{n(\mathrm{H^+})} \times \frac{k_{03} E_{32} A_{32}}{\alpha_{\mathrm{eff, H}\beta} E_{\mathrm{H}\beta} (A_{31} + A_{32})}
    \]
    可得：
    \[
        \frac{n(\mathrm{O\,III})}{n(\mathrm{H^+})} = \frac{I([\mathrm{O\,III}]5008)}{I(\mathrm{H}\beta)} \times \frac{\alpha_{\mathrm{eff, H}\beta} E_{\mathrm{H}\beta} (A_{31} + A_{32})}{k_{03} E_{32} A_{32}}
    \]
    通过查表可以得知：
    \[
        A_{31} = 6.215 \times 10^{-3},\quad A_{32} = 1.814 \times 10^{-2}
    \]
    \[
        E_{32} = \frac{hc}{5008.2\,\ang},\quad E_{\mathrm{H}\beta} = \frac{hc}{4862.7\,\ang}
    \]
    则有：
    \begin{align*}
        \frac{n(\mathrm{O\,III})}{n(\mathrm{H^+})}
        &= \frac{I([\mathrm{O\,III}]5008)}{I(\mathrm{H}\beta)} \times \frac{\alpha_{\mathrm{eff, H}\beta}}{k_{03}} \times \frac{5008.2\,\ang}{4862.7\,\ang} \times \frac{6.215 \times 10^{-3} + 1.814 \times 10^{-2}}{1.814 \times 10^{-2}} \\
        &= \frac{I([\mathrm{O\,III}]5008)}{I(\mathrm{H}\beta)} \times \frac{\alpha_{\mathrm{eff, H}\beta}}{k_{03}} \times 1.3828
    \end{align*}
    由(14.9)式得：
    \[
        \alpha_{\mathrm{eff, H}\beta} = 3.03 \times 10^{-14} T_4^{-0.874-0.058\ln{T_4}}\,\mathrm{cm^3/s} \approx 3.03 \times 10^{-14} T_4^{-0.87}\,\mathrm{cm^3/s}
    \]
    又由(18.10)式得：
    \[
        k_{03} = 8.629 \times 10^{-8} T_4^{-1/2} \frac{\Omega_{03}}{g_0} \mathrm{e}^{-E_{30}/kT}\,\mathrm{cm^3/s}
    \]
    其中$E_{30}/k = 29170\,\mathrm{K}$，$g_0 = 1$，且由表(F.2)可得：
    \[
        \Omega_{03} = 0.243 \times T_4^{0.120+0.031\ln{T_4}}
    \]
    可以近似$\Omega_{03}$与温度无关，即：
    \[
        \Omega_{03} \approx 0.243
    \]
    于是有：
    \begin{align*}
        k_{03}
        &= 8.629 \times 10^{-8} T_4^{-1/2} \frac{\Omega_{03}}{g_0} \mathrm{e}^{-E_{30}/kT}\,\mathrm{cm^3/s} \\
        &= 8.629 \times 10^{-8} T_4^{-0.5} \times 0.243 \mathrm{e}^{-29170\,\mathrm{K}/T}\,\mathrm{cm^3/s} \\
        &= 2.097 \times 10^{-8} T_4^{-0.5} \mathrm{e}^{-2.917/T_4}\,\mathrm{cm^3/s}
    \end{align*}
    则有：
    \begin{align*}
        \frac{\alpha_{\mathrm{eff, H}\beta}}{k_{03}}
        &= \frac{3.03 \times 10^{-14} T_4^{-0.874}}{2.097 \times 10^{-8} T_4^{-0.5} \mathrm{e}^{-2.917/T_4}} \\
        &= 1.445 \times 10^{-6} T_4^{-0.37} \mathrm{e}^{-2.917/T_4}
    \end{align*}
    因此：
    \begin{align*}
        \frac{n(\mathrm{O\,III})}{n(\mathrm{H^+})}
        &= \frac{I([\mathrm{O\,III}]5008)}{I(\mathrm{H}\beta)} \times \frac{\alpha_{\mathrm{eff, H}\beta}}{k_{03}} \times 1.3828 \\
        &= \frac{I([\mathrm{O\,III}]5008)}{I(\mathrm{H}\beta)} \times \ab(1.445 \times 10^{-6} T_4^{-0.37} \mathrm{e}^{-2.917/T_4}) \times 1.3828 \\
        &= 2.00 \times \frac{I([\mathrm{O\,III}]5008)}{I(\mathrm{H}\beta)} T_4^{-0.37} \mathrm{e}^{-2.917/T_4}
    \end{align*}
    即常数$C = 2.00$。

    \bigskip
    \item[18.2 (a)]
    由图(18.2)和图(18.4)可知：
    \[
        \frac{I([\mathrm{O\,III}]4364.4\,\ang)}{I([\mathrm{O\,III}]5008.2\,\ang)} = 0.003 \implies
        \begin{+cases}
            & T = 7365.18\,\mathrm{K},\quad n_e = 10^5\,\mathrm{cm^{-3}} \\
            & T = 7971.36\,\mathrm{K},\quad n_e = 10^4\,\mathrm{cm^{-3}} \\
            & T = 8055.57\,\mathrm{K},\quad n_e \leqslant 10^3\,\mathrm{cm^{-3}}
        \end{+cases}
    \]
    \[
        \frac{I([\mathrm{O\,II}]3729.8\,\ang)}{I([\mathrm{O\,II}]3727.1\,\ang)} = 1.2 \implies n_e T_4^{-1/2} = 241.524\,\mathrm{cm^{-3}}
    \]
    可见$T \approx 8000\,\mathrm{K}$，即$T_4^{-1/2} \approx 0.8^{-1/2} = 1.12$，而$n_e T_4^{-1/2} = 241.524$说明$n_e$显然小于$10^3\,\mathrm{cm^{-3}}$。

    取电子温度$T = 8055.57\,\mathrm{K}$，得电子数密度：
    \[
        n_e = \frac{241.524}{0.805557^{-1/2}} = 216.775\,\mathrm{cm^{-3}}
    \]
    综上，估计电子温度和数密度为$T = 8.06 \times 10^3\,\mathrm{K}$，$n_e = 217\,\mathrm{cm^{-3}}$。

    \medskip
    \item[18.2 (b)]
    由视星等公式：
    \[
        m - m_0 = -2.5\log\ab(\frac{I}{I_0}) + C
    \]
    可得修正后$\mathrm{O\,III}$的谱线强度比为：
    \[
        \frac{I([\mathrm{O\,III}]4364.4\,\ang)}{I([\mathrm{O\,III}]5008.2\,\ang)} = \frac{0.003}{10^{-0.4/2.5}} = 0.004
    \]
    则由图(18.2)可得电子温度为：（仍然假设$n_e \leqslant 10^3\,\mathrm{cm^{-3}}$）
    \[
        T = 8641.97\,\mathrm{K}
    \]
    得电子数密度：
    \[
        n_e = \frac{241.524}{0.864197^{-1/2}} = 224.526\,\mathrm{cm^{-3}}
    \]
    综上，估计电子温度和数密度为$T = 8.64 \times 10^3\,\mathrm{K}$，$n_e = 225\,\mathrm{cm^{-3}}$

\end{enumerate}

\end{document}